\section{Results}
\label{sec:results}

\subsection{Development and Pilot Observations}
\label{sec:pilot-results}

This subsection reports descriptive P1 development diagnostics. It is not a
final benchmark, contains one run per dataset, and supports no statistical
significance or generator-ranking claim. The corrected JSON-object transport
produced persisted parseable pools at or above the requested size: 200 records
for Abalone19, 201 for Ecoli4, and 200 for page-blocks0. The extra Ecoli4 record
is the unresolved over-generation anomaly described in
\cref{sec:reproducibility}; rates below use the recorded pool denominator.

\begin{table}[t]
  \centering
  \caption{P1 development diagnostics. Percentages use each persisted
  parseable pool as denominator. ``Unshown'' means an exact match to a TRAIN
  row not used as a prompt reference. Disagreement is B2-reject/B3-accept.}
  \label{tab:p1-pilot-diagnostics}
  \resizebox{\linewidth}{!}{%
  \begin{tabular}{@{}lrrrrrrr@{}}
    \toprule
    Dataset & Pool & Duplicate (\%) & Reference (\%) & Unshown (\%) &
      B2 accept (\%) & B3 accept (\%) & Disagree ($n$) \\
    \midrule
    Abalone19    & 200 & 28.50 & 24.50 & 0.00 & 69.00 & 100.00 & 62 \\
    Ecoli4       & 201 & 10.45 & 10.95 & 0.00 & 86.57 &  99.50 & 26 \\
    page-blocks0 & 200 &  9.50 &  1.00 & 0.00 & 87.50 & 100.00 & 25 \\
    \bottomrule
  \end{tabular}}
\end{table}


Exact duplicate rates were 28.5\% for Abalone19, 10.45\% for Ecoli4, and
9.5\% for page-blocks0. Exact prompt-reference matches accounted for 24.5\%,
10.95\%, and 1.0\% of the same pools, respectively. No dataset contained an
exact match to a TRAIN record that had not been shown as a prompt reference.
The distinction matters: these results describe exact copying categories, not
semantic similarity or near copying.

The P0-v3 to P1 scaling comparison increased the configured request from 40 to
200 records per dataset under the corrected seed protocol. Duplicate rates
rose from 17.5\% to 28.5\% for Abalone19, from 5.0\% to 10.45\% for Ecoli4,
and from 7.5\% to 9.5\% for page-blocks0. With one run at each scale, these are
descriptive observations only. They motivate repeated-seed and repeated-split
experiments; they do not estimate a scaling law.

B2 accepted 138/200 Abalone19, 174/201 Ecoli4, and 175/200 page-blocks0
records. B3 accepted every Abalone19 and page-blocks0 record and 200/201
Ecoli4 records. Thus, the current same-model, record-wise B3 baseline exhibited
little discrimination in P1; this observation does not establish that
language-model judges are generally ineffective. The disagreement report
contains 62, 26, and 25 cases where B2 rejected and B3 accepted. Of these, 57,
21, and 19 were classified as
dataset-redundancy cases, while 5, 5, and 6 were structural/domain cases. The
counts show that the policies operationalize different concepts; human labels
are required before interpreting either policy as the better record-validity
critic.

TRAIN-support summaries were successfully produced for each pool using the
representation in \cref{sec:diagnostics}. Their raw distances are not compared
across datasets because scale after feature construction, dimensionality, and
class support differ. No support threshold was selected, and support distance
did not alter an acceptance decision.

The repository also contains pilot TRTR/TSTR outputs. PAPER v0.1 does not use
them as comparative evidence: the evaluation has one split per dataset, the
minority test sets are very small, and several logistic-regression TSTR entries
are undefined in the committed reports. These artifacts remain part of the
evidence record and require protocol review before final utility claims.

\subsection{Human Ground-Truth Audit}
\label{sec:human-results}

The blinded sample contains 120 records and the committed manifest confirms
that all human-entry columns were blank when the sample was generated.
Annotation is in progress.

% TODO[HUMAN]: insert reviewer completion counts, confusion analyses, raw
% agreement, Cohen's kappa, adjudication outcomes, and uncertainty intervals.
\papertodo{HUMAN}{Results require completed blinded annotation and adjudication.}

\subsection{Final Benchmark}
\label{sec:final-benchmark-results}

No final benchmark has been run. Dataset inclusion, repeated-split design,
uncertainty intervals, and statistical tests must be frozen before this
subsection is populated.

% TODO[RESULTS]: replace all planned table cells only from versioned final
% benchmark artifacts.
\papertodo{RESULTS}{Final benchmark and statistical analysis are pending.}

\subsection{Ablations}
\label{sec:results-ablations}

No proposed M2+ component is implemented, so no ablation outcome is available.
The planned factorial questions are specified in \cref{sec:ablations}; detailed
future-only table shells are retained in the non-rendered manuscript notes.

\begin{table}[t]
  \centering
  \caption{Pending human critic benchmark. All entries require completed
  annotation and adjudication.}
  \label{tab:future-human-benchmark}
  \begin{tabular}{@{}lrrrr@{}}
    \toprule
    Critic & Precision & Recall & F1 & Agreement \\
    \midrule
    B2 deterministic & -- & -- & -- & -- \\
    B3 single judge & -- & -- & -- & -- \\
    Proposed fused evidence & -- & -- & -- & -- \\
    \bottomrule
  \end{tabular}
\end{table}

\begin{table}[t]
  \centering
  \caption{Planned final B1/B2/B3/\rcqasyn{} comparison; \rcqasyn{} is not
  implemented.}
  \label{tab:future-baselines}
  \begin{tabular}{@{}lrrrr@{}}
    \toprule
    System & Human validity & Coverage & TSTR & Cost \\
    \midrule
    B1 raw & -- & -- & -- & -- \\
    B2 deterministic & -- & -- & -- & -- \\
    B3 single judge & -- & -- & -- & -- \\
    Proposed \rcqasyn{} & -- & -- & -- & -- \\
    \bottomrule
  \end{tabular}
\end{table}

